% ============================================================
%  Real-time engine
%  Streaming architecture: NTRIP transport, RTCM/MSM decode,
%  observation/SSR time synchronisation, per-epoch RT update.
%
%  Sources: src/core/realtime/{RealTimeEngine,NtripClient,
%  RtcmStreamReader}.cpp, src/core/rtcm/{RtcmParser,RtcmMsmDecoder,
%  RtcmSsrDecoder,RtcmNavDecoder}.cpp, src/core/sbf/{SbfStreamReader,
%  SbfParser,Meas3Decoder,SbfSignalMap,SubframeAssembler}.cpp,
%  src/core/app/ApplicationRealtime.cpp, src/core/math/{TimeUtils,
%  Coordinates}.cpp, src/core/rinex/RinexParser.cpp,
%  src/core/products/{SsrCorrectionStore,LiveEphemerisStore}.cpp.
% ============================================================
\documentclass[main.tex]{subfiles}
\begin{document}
\ifSubfilesClassLoaded{\begin{center}\small\itshape Part of the GnssProcessor Mathematical Reference \textbullet\ \today\end{center}}{}

\section{Real-time engine}\label{sec:rt}

The real-time path reuses the estimator, measurement model and state
manager of post-processing and reorganises the control flow around a
\emph{streaming} contract.  Complete RTCM\,3 (or SBF) frames arrive one
at a time from a file or a network socket and are decoded into
observations, corrections and broadcast ephemeris.  A buffered epoch
passes through the \dcref{sec:kf}{Kalman-filter} predict/update cycle as
soon as its boundary is detected.  Post-processing iterates a closed
list of epochs in time order with the full orbit and clock product set
on disk; the engine sees each epoch once, receives corrections that may
lead or lag the observations, and cannot look ahead.  The entry point is
\texttt{Application::runRealTime} (\texttt{src/core/app/Application\-Realtime.cpp}),
which constructs a \texttt{RealTimeEngine}
(\texttt{src/core/realtime/RealTimeEngine.cpp}) and drives it from a
single-threaded multi-source poll loop.  The engine takes its input
through the frame feeds \texttt{feedRtcmFrame} and \texttt{feedSbfFrame};
at stream end \texttt{flush} processes whatever epochs are still
buffered.  This section covers the parts specific to that path.  The estimation mathematics, the observation
equations and the SSR correction model are documented elsewhere and are
only cross-referenced here.

\subsection{Streaming control flow}\label{sec:rt-flow}

The poll loop holds one \texttt{RealtimeStream} per logical source and
services the sources in turn without blocking on any one of them:
\begin{itemize}[leftmargin=*, itemsep=0pt]
\item a \emph{primary} stream carrying rover observations (RTCM\,3 MSM
      or Septentrio SBF, from a file replay, an NTRIP mountpoint, or a
      raw TCP socket);
\item an optional \emph{corrections} stream (typically an IGS-RTS NTRIP
      mountpoint carrying SSR);
\item an optional \emph{ephemeris} stream (broadcast nav for an
      obs-only station that does not co-broadcast 1019/1020/\ldots).
\end{itemize}
Each iteration drains a bounded burst from every active source through
its framer and dispatches the complete frames to the engine
(\texttt{pumpFrames}).  The primary drain is capped at
\SI{1}{\mebi\byte} per iteration, and the corrections and ephemeris
streams are read one \SI{4}{\kibi\byte} chunk per iteration, so a
high-rate observation stream cannot starve the sparse SSR and navigation
streams.  When a live network-product tail is configured, the loop also
polls it at about \SI{5}{\hertz} and wakes the obs/SSR hold buffer of
\Cref{sec:rt-sync} when new products arrive.  The loop exits when every
source is inactive (file EOF or client stopped), when a stop is
requested (Ctrl+C or the console-close handler), or when the
\texttt{processing.max\_duration\_s} data-time budget latches.  On exit
the engine flushes the final buffered epoch before teardown, so no
observations or trace rows are lost.

A socket source counts as \emph{active} while its client is running,
whether or not it is connected at that moment.  A transient drop or the
pre-connect window at startup therefore does not end the session; the
client reconnects in the background.

\subsection{NTRIP transport}\label{sec:rt-ntrip}

\paragraph{Caster connection.}
\texttt{NtripClient} (\texttt{src/core/realtime/NtripClient.cpp}) is an
NTRIP\,v2 client built on libcurl.  A background thread holds the HTTP
GET open and appends the received bytes to a mutex-guarded buffer with a
hard cap of \SI{4}{\mebi\byte}.  When the cap is reached the buffer
drops its oldest bytes.  Consumers read through a blocking
\texttt{read} or a timeout-bounded \texttt{tryRead}; the multi-stream
loop uses the latter, so one stream cannot block the others.  The
request carries the NTRIP\,v2 headers (\texttt{Ntrip-Version:
Ntrip/2.0} and a \texttt{User-Agent}), and the mountpoint is the URL
path.  HTTP Basic authentication (\texttt{CURLAUTH\_BASIC}) is supplied
when a username is configured.  An optional NMEA GGA sentence for VRS
mountpoints is sent as the NTRIP\,v2 \texttt{Ntrip-GGA} request header on
the GET; the request is never turned into a POST.

\paragraph{Transport security.}
The transport defaults to plain HTTP, which serves plain-HTTP casters
such as IGS-IP (\texttt{igs-ip.net}) unchanged.  A \texttt{use\_tls}
flag selects the \texttt{https} URL scheme and sets
\texttt{CURLUSESSL\_ALL}, as required for port-443 casters such as CDDIS
that offer no plain-HTTP endpoint.  Certificate verification
(\texttt{SSL\_VERIFYPEER}/\texttt{VERIFYHOST}) stays on unless a
testing-only \texttt{tls\_insecure} override disables it for self-signed
or mis-chained casters.  Both flags default off, so a plain-HTTP run is
unaffected.

\paragraph{Reconnection.}
The receive loop reconnects automatically.  After each transfer ends it
classifies and logs the outcome: a transport error, a non-200 HTTP
status (with hints for 401, 403 and 404), a connection that delivered
zero bytes (an idle mountpoint, or a VRS mountpoint that needs a GGA
position), or a clean stream end.  It then always waits
\texttt{reconnect\_delay\_s} plus a uniform 0--2000\,ms jitter before
retrying, so a misconfigured mountpoint appears in the log as an error
and several streams dropped by one caster outage do not reconnect in
lockstep.  An NTRIP\,v1 caster answers ``ICY 200 OK'', which libcurl
reports as status code 0; the loop treats code 0 as success.

\subsection{RTCM\,3 framing and CRC}\label{sec:rt-framing}

\texttt{RtcmStreamReader} (\texttt{src/core/realtime/RtcmStreamReader.cpp})
turns an arbitrary byte stream into discrete RTCM\,3 frames and so
separates transport from parsing.  An RTCM\,3 frame is
\begin{equation}
\underbrace{\texttt{0xD3}}_{\text{preamble}}\;
\underbrace{[\,6\text{ reserved}\,|\,10\text{-bit length }L\,]}_{\text{2 bytes}}\;
\underbrace{[\,\text{payload, }L\text{ bytes}\,]}_{}\;
\underbrace{[\,\text{CRC-24Q}\,]}_{3\text{ bytes}},
\label{eq:rt-frame}
\end{equation}
so the total frame length is $3 + L + 3$.  \texttt{nextFrame} scans for
the \texttt{0xD3} preamble (discarding other bytes), reads the 10-bit
payload length (a length above 1023 bytes is rejected as a false
preamble), waits until the full $3+L+3$ bytes are buffered, and verifies
the trailing 24-bit Qualcomm CRC (CRC-24Q, table-driven) computed over
header$+$payload.  On a CRC mismatch the reader advances past the
candidate preamble byte and resumes scanning, so a corrupt or misframed
byte cannot stall it.  Only CRC-valid frames are emitted, and the
downstream parser (\texttt{RtcmParser}) relies on frame integrity.  The
reader counts extracted frames and discarded bytes for stream-health
diagnostics.

\paragraph{Septentrio SBF.}
A parallel binary path decodes Septentrio SBF frames into the same
observation and ephemeris outputs (\texttt{src/core/sbf/}).
\texttt{SbfStreamReader} performs the \texttt{\$@} (\texttt{0x24\,0x40})
sync and the CRC-16/XMODEM check, and \texttt{SbfParser} dispatches
blocks by ID and decodes the \texttt{MeasEpoch} observations.
\texttt{Meas3Decoder} unpacks the compact, differentially encoded
\texttt{Meas3Ranges} observations; GLONASS FDMA slots are resolved
through \texttt{SbfSignalMap}.  \texttt{SubframeAssembler} rebuilds
broadcast ephemeris from raw-nav subframes.  This section covers only the
framing and the feed seam (\Cref{sec:rt-flow,sec:rt-apriori}).  The
\texttt{Meas3Ranges} differential range reconstruction is the direct
analogue of \eqref{eq:rt-msm-code} and \eqref{eq:rt-msm-phase} below and
the one part of the SBF path that is algebra rather than ICD
bookkeeping.  It is \emph{not} yet derived in this reference, and no
other document derives it: this is a known gap.  The remaining SBF detail
(\texttt{SubframeAssembler} ICD scaling, the \texttt{SbfSignalMap}
tables, the framing constants) belongs in the Doxygen headers and is out
of scope here.

The top-level \texttt{RtcmParser} (\texttt{src/core/rtcm/RtcmParser.cpp})
reads the 12-bit message type DF002 from the first 12 bits of the
payload (the frame header carries only the length) and dispatches on it:
1005/1006 station coordinates; 1033 receiver and antenna descriptors and
1230 GLONASS code-phase biases (both logged, neither consumed); 4076
IGS-SSR; legacy SSR (1057--1068, 1240--1270); broadcast ephemeris
(1019/1020/1042/1044/1046); and MSM4, MSM5 and MSM7 observations.  The
week that promotes a truncated epoch field to a full \texttt{GnssTime}
comes from the parser's per-stream time tracker rather than from the
host clock at each frame (\Cref{sec:rt-time}).  One parser instance
serves the primary, corrections and ephemeris streams, so all three share
one GPS timeline.

\subsection{MSM observation decode}\label{sec:rt-msm}

A single body (\texttt{decodeMsmImpl},
\texttt{src/core/rtcm/RtcmMsmDecoder.cpp}) decodes every MSM observation
message.  It is parameterised over an \texttt{MsmVariant} descriptor that
supplies the per-field bit widths, scale factors, sentinels and
capability flags.  Three variants are instantiated: MSM4, MSM5 and MSM7.
The variant mechanism can describe MSM6, but no MSM6 entry exists, so
MSM6 is not decoded.  The header carries a 64-bit satellite mask, a
32-bit signal mask and an $N_{\text{sat}}\times N_{\text{sig}}$ cell
mask; the number of cells is
$N_{\text{cell}} = \operatorname{popcount}(\text{cell mask})$.  Each
field block is packed contiguously across all satellites or all cells, in
mask order, and is decoded into parallel vectors before recombination.

\paragraph{Range reconstruction.}
For each satellite the rough range in milliseconds is the sum of the
8-bit integer part (DF397) and the 10-bit modulo part
($2^{-10}\,\text{ms}$, DF398).  Each cell adds a signed fine term with a
variant-specific scale.  With $c$ the speed of light, the pseudorange and
carrier phase of a cell are
\begin{align}
\code &= c\cdot 10^{-3}\Big(
  n_{\text{int}} + n_{\text{mod}}\,2^{-10}
  + n_{\text{fine}}^{\code}\,\sigma_{\code}\Big),
\label{eq:rt-msm-code}\\
\phaseo_{[\text{m}]} &= c\cdot 10^{-3}\Big(
  n_{\text{int}} + n_{\text{mod}}\,2^{-10}
  + n_{\text{fine}}^{\phaseo}\,\sigma_{\phaseo}\Big),
\qquad
\phaseo_{[\text{cyc}]} = \phaseo_{[\text{m}]} / \lambda,
\label{eq:rt-msm-phase}
\end{align}
with fine scales $\sigma_{\code}=2^{-24}$, $\sigma_{\phaseo}=2^{-29}$
for MSM4 (and MSM5) and $\sigma_{\code}=2^{-29}$, $\sigma_{\phaseo}=2^{-31}$ for
MSM7 (all in ms).  Cells holding the RTCM-defined invalid sentinel are
skipped.  Carrier phase is emitted in cycles, the RINEX-3 convention the
observation model consumes (\dcref{sec:obs}{the observation-equations
section}); when the wavelength $\lambda$ cannot be resolved, the phase
is stored in metres under the same observation code.  The decoder also
produces CNR (DF403 in whole dB-Hz for MSM4/5, DF408 in 1/16 dB-Hz for
MSM7) and, for MSM5/7, a Doppler value from the rough$+$fine
phase-range-rate (the negated rate over the wavelength).  The
half-cycle-ambiguity bit
(DF420) sets the LLI half-cycle flag.

\paragraph{DF419 extended-satellite-info gating.}
The 4-bit DF419 ``extended satellite information'' field determines the
framing of the per-satellite block.  In MSM5 and MSM7 it is present in
the bitstream for every constellation, but only GLONASS assigns it a
meaning (the FDMA frequency-channel number).  The decoder therefore
gates two operations separately:
\begin{itemize}[leftmargin=*, itemsep=0pt]
\item \emph{bitstream consumption} is gated on the MSM variant
      (\texttt{honour\_ext\_sat\_info}, true for MSM5/MSM7, false for
      MSM4/MSM6): when set, $4\,N_{\text{sat}}$ bits are read for all
      systems;
\item \emph{interpretation} is gated separately on the system
      descriptor; only GLONASS maps the value to a channel.
\end{itemize}
Reading DF419 only for GLONASS (or skipping it in MSM7) under-consumes
the per-satellite block by $4\,N_{\text{sat}}$ bits, so the per-cell
fine pseudoranges are read from the wrong bit offset.  The rough range
survives but the fine terms are biased, which produces an L1/L2 split of
order hundreds of metres with the wrong sign on the geometry-free
combination.  These observations look plausible and pass an over-read
check.  Gating on the variant is the fix for the GPS L1/L2 geometry-free
sign anomaly, and it matches RTCM\,10403.3 and the cssrlib reference.  As a
second guard, the parser's \texttt{warnIfUnderConsumed} warns whenever a
decode leaves $\ge 8$ payload bits unread, which detects this class of
misalignment where the bit reader's over-read flag does not.

\subsection{Broadcast nav, SSR and station coordinates}\label{sec:rt-other}

\paragraph{Broadcast ephemeris.}
Messages 1019 (GPS LNAV), 1044 (QZSS), 1046 (Galileo I/NAV), 1042
(BeiDou D1) and 1020 (GLONASS) are decoded by \texttt{RtcmNavDecoder}
into a common \texttt{BroadcastEphemeris} and pushed into the
process-wide \texttt{LiveEphemerisStore}, from which the geometry layer
resolves satellite orbits when no SP3/CLK file is loaded.  The field
order differs by message (QZSS, for example, leads with its clock
block).  GLONASS carries a sign-magnitude PZ-90 state vector, which is
integrated separately.

\paragraph{SSR corrections.}
\texttt{RtcmSsrDecoder} decodes both the legacy SSR messages and the
IGS-SSR 4076 container and pushes the records into the shared
\texttt{SsrCorrectionStore} with their SSR header (provider epoch and
IOD-SSR), so staleness and IOD diagnostics refer to the provider
solution rather than to the time of receipt.  The store holds orbit,
clock, combined, code-bias, phase-bias, URA and high-rate-clock records
per satellite.  The orbit and clock query paths read the store directly,
so no per-epoch wiring is needed.  The correction mathematics (the
radial/along/cross orbit projection, the clock and bias application, the
ionosphere model) is documented in \dcref{sec:ssr}{the SSR-corrections
section} and is not repeated here.  Two reactions to the stream are
specific to the real-time path:
\begin{itemize}[leftmargin=*, itemsep=0pt]
\item an incremented IGS-SSR phase-bias \emph{discontinuity counter}
      for a signal breaks the integer property of that signal's
      ambiguity.  By default the engine resets the affected ambiguity
      before the same epoch's measurement update
      (\texttt{StateManager::noteSsrDiscontinuity} detects the counter
      change).  Two opt-in knobs replace the reset: the declared-step
      carry \texttt{processing.ssr\_phase\_bias\_disc\_carry} (card CC)
      and the step translation
      \texttt{processing.ssr\_phase\_bias\_step\_translate} (card PBT,
      which also covers undeclared steps and takes precedence).  Both shift the ambiguity by the
      phase-bias step instead of re-floating it;
\item every SSR push re-attempts a held observation epoch
      (\Cref{sec:rt-sync}), because the correction horizon has just
      advanced.
\end{itemize}

\paragraph{Station coordinates (1005/1006).}
\texttt{decodeStationCoords} reads the antenna reference point (ARP) as
three signed 38-bit ECEF components in \SI{0.1}{\milli\metre} units;
1006 adds a 16-bit antenna height in the same unit.  A reference
station that streams its own coordinates thereby supplies an a-priori
seed (\Cref{sec:rt-apriori}).

\subsection{Time-system handling}\label{sec:rt-time}

Streaming messages carry truncated, system-specific epoch fields, which
are promoted to a common GPS-time \texttt{GnssTime}.  Each parser
instance keeps a tracker, in continuous GPS seconds, of the last decoded
stream epoch.  It promotes every mod-week or time-of-day field by
\emph{nearest continuation} to that tracker: the week (or day) that
minimises $|\Delta t|$, as in RTKLIB's \texttt{adjweek}
(\texttt{RtcmParser::anchorGpsSeconds}).  On this path the host clock
(\texttt{TimeUtils::nowContinuousGpsSeconds}, POSIX UTC plus the leap
offset) is read only at cold start, on the first time-bearing frame
after construction, so a skewed host clock or a race at a week rollover
cannot move an established stream to another week.  Each successful
MSM or SSR epoch advances the tracker (\texttt{updateTimeAnchor}).  The
Keplerian broadcast-nav decoders read the tracker for their week but do
not advance it.  The anchor below is the tracker's current value.  The
1020 GLONASS ephemeris \textsc{toe} is dated from the same tracker; a
1020 that arrives before any MSM or SSR frame has set it is held and
dated by the first such frame (see the GLONASS item below).
\begin{itemize}[leftmargin=*, itemsep=0pt]
\item \textbf{Flat ms-of-week systems} (GPS/GAL/QZS/SBAS/BDS MSM):
      the 30-bit epoch is milliseconds of week.  BeiDou adds the
      constant \SI{14}{\second} BDT$\to$GPS offset and re-wraps the
      result into the week.
\item \textbf{GLONASS} carries its 30-bit MSM epoch as a 3-bit
      day-of-week followed by a 27-bit time-of-day in ms.  Reading only
      27 bits misaligns the satellite mask and corrupts the whole
      GLONASS decode.  \texttt{TimeUtils::glonassTodToGpsTime} promotes
      the ms-of-day: it places the reported time-of-day in the
      GLONASS\,$=$\,UTC(SU)$+$\SI{3}{\hour} (``Moscow'') day nearest the
      stream-time tracker and converts back to GPS time, re-adding the
      leap seconds:
\begin{equation}
t_{\text{GPS}} =
  \big(\,t_{\text{day}}^{0} + \tilde{t}_{\text{TOD}}\,\big)
  - \Delta_{\text{Moscow}} + \Delta_{\text{leap}},
\label{eq:rt-glo-time}
\end{equation}
      where $t_{\text{day}}^{0}$ is the Moscow-day midnight nearest the
      anchor and $\tilde{t}_{\text{TOD}}$ is the reported time-of-day
      placed into that day by resolving the midnight wrap toward the
      nearest match ($\pm\SI{12}{\hour}$).  The 3-bit day-of-week is
      read but not used, since the anchor resolves the day.  The same
      conversion promotes any GLONASS SSR epoch against the tracker.  It
      also promotes the 1020 ephemeris \texttt{toe} against the same
      tracker, so a file replay of historical data keeps each GLONASS
      ephemeris on the day of its observations.  A 1020 that arrives
      before any MSM or SSR frame has set the tracker has no stream
      epoch to be dated from.  It is held (at most 48 records) and dated
      by the first such frame; an input that carries navigation messages
      only fills the hold, which is then dated from the host clock.
\item \textbf{GLONASS FDMA wavelength.}  GLONASS L1/L2 are
      frequency-division: each slot transmits on
      $f = f_{\text{base}} + k\,\Delta f$ with channel
      $k \in [-7,+6]$, so the carrier wavelength is per satellite.  The
      decoder maps DF419 to the channel as $k = v - 7$ for the valid
      range and uses it in \Cref{eq:rt-msm-phase}.  The channel is also
      stored on the observation, so the slot map and the MW combinations
      can resolve frequency.  When DF419 is absent (MSM4/MSM6) or holds a
      reserved value, the channel is marked unknown and the phase is
      stored in metres; a default of zero would decode to a false
      $k=-7$.
\end{itemize}
Each decoded record carries its own content epoch
(\texttt{epoch\_time}).  The synchronisation logic of
\Cref{sec:rt-sync} compares against that epoch, not against the time of
arrival.

\subsection{Observation/SSR synchronisation}\label{sec:rt-sync}

PPP-RTK needs observations time-\emph{aligned} with the corrections
rather than with the host clock.  In a live session the observation
stream commonly runs ahead of a lagging SSR stream.  An observation
epoch that no satellite's SSR can cover is processed entirely at
broadcast grade (and, with \texttt{drop\_broadcast\_grade}, loses every
satellite).  The opt-in \texttt{processing.enable\_obs\_ssr\_sync} hold
buffer prevents this.  With $t_{\text{obs}}$ the observation epoch, the
SSR coverage horizon is the freshest stored SSR content epoch
(\texttt{SsrCorrectionStore::newestSsrEpoch}) plus the maximum SSR age:
\begin{equation}
t_{\text{horizon}} =
  t_{\text{SSR}}^{\text{newest}} + \tau_{\text{max}},
\qquad
\tau_{\text{max}} = \texttt{ssr\_max\_age\_s}.
\label{eq:rt-horizon}
\end{equation}
The decision is:
\begin{itemize}[leftmargin=*, itemsep=0pt]
\item $t_{\text{obs}} \le t_{\text{horizon}}$: SSR covers the epoch.  The
      epoch is processed normally and any hold is cleared;
\item $t_{\text{obs}} > t_{\text{horizon}}$: the epoch is
      \emph{deferred}.  The hold buffer keeps only the latest such epoch
      (skip-to-latest semantics), and \texttt{drainHeldObs} re-attempts
      it as soon as an SSR push (or a new network-product epoch,
      through \texttt{notifyProductsAdvanced}) advances the horizon.
\end{itemize}
A liveness backstop keeps a dead SSR feed from starving the filter.  If
a hold persists past \texttt{obs\_ssr\_sync\_max\_hold\_s} of host-clock
time, the engine latches a degraded ``outage'' state and passes every
epoch through (at broadcast grade, or dropped by the existing path)
until SSR covers the observations again.  The held epoch is drained with
skip-to-latest \emph{bypassed}, so the SSR-aligned epoch is not
discarded.  The final \texttt{flush} bypasses both the hold and
skip-to-latest, so the last epoch is always processed.  The feature is
inert when disabled, before the first epoch, and in any run without a
live SSR orbit and clock source.

A separate, coarser \texttt{enable\_obs\_skip\_to\_latest} drops
observation epochs whose lag behind the \emph{host-clock} GPS time
exceeds \texttt{obs\_skip\_max\_lag\_s} (default \SI{5}{\second}), so a
backlogged stream
fast-forwards to the freshest epoch.  The skip returns before any
continuity state is touched, so the next processed epoch measures its
time step from the last accepted one and trips the gap path
(\Cref{sec:rt-epoch}) when the skipped span exceeds the gap tolerance.

\subsection{Ephemeris management}\label{sec:rt-eph}

\texttt{LiveEphemerisStore} keeps, per satellite, the current record and
one previous record (the current record is demoted to previous when the
IODE changes), and evaluates the orbit and clock against the appropriate
one.  Two policies govern it in real time.
\begin{itemize}[leftmargin=*, itemsep=0pt]
\item \textbf{Freshest-\textsc{toe} selection.}  An incoming record is
      rejected when its \texttt{toe}, expressed in continuous GPS
      seconds (\texttt{toeGpsContinuous}, which reconciles the BDT and
      GLONASS time scales), is older than the held record's by more than
      \SI{1}{\second}.  A stale or re-ordered broadcast therefore never
      overwrites a newer one.  A record from the IGS stream is also
      protected from a receiver-decoded record at equal \texttt{toe}:
      the IGS-RTS SSR deltas are computed against the IGS broadcast, and
      a receiver ephemeris with different fit parameters would make the
      pairing inconsistent.  Receiver nav replaces an IGS-stream record
      only when its reference epoch is at least \SI{1800}{\second} newer
      (the GLONASS nominal \texttt{toe} step), which indicates that the
      IGS stream has stalled on that satellite.
\item \textbf{\textsc{toe}-aware and IOD-aware selection.}  For orbit
      and clock evaluation the store picks the record whose
      \texttt{toe} is nearest the query epoch (\texttt{pick}).  For an
      SSR-corrected satellite it first selects the record the SSR delta
      was computed for (\texttt{pickForSsr}, which matches the SSR IODE
      against the current and the retained previous record).  Across an
      ephemeris cutover the receiver can be one IOD ahead of the SSR
      stream; without this re-reference the IOD mismatch would drop the
      satellite to broadcast grade.
\end{itemize}

\subsection{A-priori seeding from RTCM 1005/1006}\label{sec:rt-apriori}

A \emph{live} SBF/RTCM observation stream carries no RINEX \texttt{APPROX
POSITION} record (a file-fed source does; see the third item below).
Without an a-priori position the engine falls back to a coarse
Earth-centred seed at the geographic pole.  \texttt{seedInitialState}
calls \texttt{PppModel::initializeState} and then applies the fallback,
since a state write before that call would be overwritten.  The model it
calls is built by \texttt{initProcessingComponents}, so
\texttt{seedInitialState} must follow it in the construction order.  From there almost every
satellite sits below the elevation mask, the first
\texttt{computeModels} yields no usable satellite, WLS never runs and
the filter cannot bootstrap.  The engine therefore seeds the a-priori
position from the stream itself until the Kalman filter owns the
position:
\begin{itemize}[leftmargin=*, itemsep=0pt]
\item from a \emph{PVTGeodetic} solution in an SBF stream
      (\texttt{feedSbfFrame});
\item from an RTCM \emph{1005/1006} station ARP
      (\texttt{feedRtcmFrame}), the RTCM analogue of the PVT seed.  This
      makes \texttt{receiver.reference\_position} optional for any
      station that streams 1005/1006.  An optional blunder gate
      (\texttt{realtime.arp\_blunder\_gate\_m}, RTA-12(3); 0, the
      default, disables it) rejects a broadcast ARP that disagrees with
      an existing a-priori by more than that distance, so a wrong
      mountpoint or stale 1005 coordinates cannot move the seed;
\item from the RINEX header of a \emph{file-fed} source, reduced from the
      marker to the antenna reference point through
      \texttt{ANTENNA: DELTA H/E/N}
      (\texttt{RinexParser::getArpPositionXyz}, which calls
      \texttt{Coordinates::markerToArp}).  The position state, and the
      ATX PCO/PCV applied to it, refer to the ARP rather than the
      monument, so seeding on \texttt{APPROX POSITION XYZ} unreduced
      would start the filter one eccentricity away.  The reduction is an
      exact no-op for a zero or absent eccentricity, so a live SBF/RTCM
      stream is unaffected; it matters only for the file-fed closed-loop
      consumer, whose header carries a real eccentricity.  The reduction
      is not applied twice: \texttt{PppModel::initializeState} reduces
      from the \emph{raw} marker, and only while the state position is
      still zero or non-finite.
\end{itemize}
A configured \texttt{reference\_position} always takes precedence.  The
seed is applied only while \texttt{first\_epoch\_} is set.  The startup
gate keeps that flag set through warmup, so a 1005/1006 message that
arrives a few seconds into the stream (these messages are sent less
often than the \SI{1}{\hertz} observations) still seeds before the first
estimation epoch.  The flag clears in the first epoch that passes the
startup gate, after measurement conditioning and before the predict
(\Cref{sec:rt-epoch}).  From then on the filter owns the position, and a
decoded ARP or PVT is only logged for comparison.

\subsection{The per-epoch real-time update}\label{sec:rt-epoch}

\texttt{processEpoch} runs the predict/update pipeline of
post-processing, with streaming-specific guards in front and the
post-processing-only stages removed.  It is a thin ordered driver.  Each
step is a private member function (a \emph{seam}) in one of the files
\texttt{src/core/realtime/RealTimeEngine*.cpp}, grouped by concern:
admission (\texttt{RealTimeEngineAdmit.cpp}), measurement-set
preparation (\texttt{RealTimeEnginePrepare.cpp}), slip QC
(\texttt{RealTimeEngineSlipQc.cpp}), SSR datum-break reactions
(\texttt{RealTimeEngineSsrBreaks.cpp}) and estimation
(\texttt{RealTimeEngineEstimate.cpp}).  Seams exchange per-epoch values
through \texttt{RtEpochContext}, which holds nothing beyond the epoch; a
value that must survive to the next epoch is a member of the engine.
The four admission gates are the only seams that may end an epoch early.
The driver itself also ends an epoch when no observation of an active
system remains or when \texttt{computeModels} fails; after the startup
gate every step runs unconditionally.  An optional staged filter-state
snapshot (warm resume) is restored before the gates, so the first time
step spans the actual outage.  In order:
\begin{enumerate}[leftmargin=*, itemsep=0pt]
\item \textbf{Lag, sync and duration gates.}  The host-clock backlog
      warning, the obs$\leftrightarrow$SSR hold of \Cref{sec:rt-sync}
      and the skip-to-latest drop form the first gate
      (\texttt{admitBacklogAndSync}); the \texttt{max\_duration\_s} latch
      is the second (\texttt{gateMaxDuration}).  Both run before any
      continuity state is touched.
\item \textbf{Step and gap detection} (the third gate,
      \texttt{stepTimeAndDetectGap}).
      $\dt = t_k - t_{k-1}$, with a week-rollover correction; the first
      epoch assumes \SI{1}{\second}.  A non-advancing $\dt \le 0$ is a
      duplicate (a second co-timestamped SBF block or a retransmit) and
      is skipped with the state untouched.  Folding it in would
      double-count the measurements and pass $\dt = 0$ to the
      process-noise build, which throws.  A $\dt$ above the gap
      tolerance flags a gap.  The tolerance is
      \texttt{processing.data\_gap\_tolerance\_s} when set.  When unset
      it is 1.5 times the median spacing of the epochs accepted so far
      (\texttt{EpochCadence}), as in post-processing, but never below
      \SI{1}{\second}: that is also the value until a first spacing has
      been measured, and a stream of \SI{1.5}{\hertz} or faster keeps it.
\item \textbf{Models and startup gate} (the fourth gate,
      \texttt{gateCorrectionsReady}).  Observations
      are filtered to the active systems and run through
      \texttt{computeModels}.  When
      \texttt{defer\_until\_corrections\_ready} is armed, estimation (no
      predict, no update) is deferred until enough cm-grade
      broadcast$+$SSR satellites are present: 5 in total and 4 per
      active constellation by default, the per-constellation floor
      reduced to what the constellation currently fields but never
      below two.  A timeout forces a degraded start.  The filter then
      initialises code and phase together on a clean state instead of
      cold-starting on metre-grade broadcast clocks.
\item \textbf{Measurement-set conditioning.}  After two observability
      sweeps (the SSR coverage log and the SSR consumer trace),
      \texttt{dropBroadcastGrade} removes broadcast-grade satellites
      when the startup gate's \texttt{drop\_broadcast\_grade} is set.  In
      modes with a receiver-clock state, \texttt{seedClockFromWls}
      replaces the filter clock with a WLS clock estimate when the two
      differ by at least \SI{5e-5}{\second} (\texttt{kJumpThresh},
      \SI{0.05}{\milli\second}, about \SI{15}{\kilo\metre}), after a
      gap, or on the first epoch.  Under BSSD (no receiver-clock state)
      \texttt{regaugeBssd} subtracts the
      median of the epoch's code innovations, one receiver clock common
      to all satellites, from every code and phase observable.  This
      common-mode shift leaves the single differences
      exactly unchanged and keeps the per-satellite bias states at DCB
      scale instead of absorbing the full unsteered receiver clock.  The
      clock seed and the BSSD re-gauge are mutually exclusive: the first
      requires a receiver-clock state, the second its absence.  The
      first-epoch flag (\texttt{first\_epoch\_}) clears at the end of
      this step.
\item \textbf{Predict.}  \texttt{predictAided} over $(t_{k-1}, t_k]$,
      the aiding path shared with post-processing; for a GNSS-only
      stream it delegates byte-identically to \texttt{predict}$(\dt)$ of
      \dcref{sec:kf}{the estimation-engine section}.
\item \textbf{Cycle-slip detection and repair}
      (\texttt{detectAndRepairSlips}).  The MW detector, the
      optional TurboEdit MW$+$GF joint repair and the soft-LLI gap
      bridge run before the update, with the detectors and repair
      machinery of \dcref{sec:slip}{the cycle-slip section}.  They are
      enabled under \texttt{estimation.cycle\_slip}.  TurboEdit is for
      high-rate data only; its GF-arc gap tolerance is 1.5 times the
      configured stream rate (\SI{1}{\hertz} if none is set), because a
      live stream has no RINEX interval header.
\item \textbf{SSR datum-break reactions.}  These react to the producer
      moving its datum and do nothing without SSR content.  The order is
      fixed.  A change of the SSR solution-generation id
      (\texttt{reactToGenerationChange}) resets every ambiguity and
      re-inflates the receiver code-bias states; it runs first because it
      subsumes the per-signal reactions.  \texttt{logProducerStatus} logs
      a producer status transition.  \texttt{resetOnPhaseBiasDisc}
      resets (or, with cards CC or PBT, shifts) the ambiguity of a signal
      whose phase bias shows a discontinuity or a freshness flip, as
      described for the SSR corrections above.
      \texttt{reinflateOnCodeBiasDisc} re-inflates the receiver
      code-bias states on a code-bias datum re-reference, and
      \texttt{reinflateOnClockDisc} re-inflates the affected
      inter-system bias states on a clock datum re-reference.  Last, the
      opt-in declared-level-step reaction
      \texttt{reactToDeclaredLevelSteps\_} (card CLV2,
      \texttt{processing.ssr\_level\_step\_inflate}, default off)
      compares each satellite's producer-declared clock and code-bias
      levels with their previous values.  When the
      resulting step in a code signal's level exceeds a
      \SI{0.1}{\milli\metre} dead band, it inflates the covariance along
      the direction that moves the modelled code but not the modelled
      phase (the satellite's slant ionosphere together with its
      ambiguities, each ambiguity weighted by the squared ratio of the
      GPS L1 frequency to its signal frequency).  The added standard
      deviation is the step size times a configured scale.  A satellite
      without an ionosphere state or an ambiguity is refused.
\item \textbf{Measurement update} (\texttt{runMeasurementUpdate}).  The
      outlier and
      forced-pivot-handover rebuild loop of the post-processing driver
      (shared via \texttt{EpochUpdateSupport}): rebuild on
      \texttt{NeedsMeasurementRebuild}, forward each constellation's
      forced pivot handover, re-inflate reset-satellite states, clear
      the MW/TurboEdit arc history of KF-reset signals, and remove one
      outlier per pass.  Phase signals held by the soft-LLI bridge are
      excluded from the update.
\item \textbf{Pseudo-observation constraints.}  After the update and
      before AR, so the reduced covariance reaches the resolver.  An
      optional soft a-priori ZWD pseudo-observation
      (\texttt{applyZwdConstraint}, \texttt{estimation.zwd\_constraint})
      anchors the BSSD [up,\,ZWD] near-null subspace.  When the engine
      consumes a network ionosphere product,
      \texttt{applyNetworkIonoPrior} constrains each satellite's
      ionosphere state with a slant-ionosphere prior from the same
      epoch, with an error block common to each constellation.
\item \textbf{Ambiguity resolution} (\texttt{resolveAmbiguities}).
      After a successful float update,
      MLAMBDA AR runs on the committed state with the configuration
      surface of post-processing and is fed the MW-resolved wide-lane
      integers (\dcref{sec:ar}{the ambiguity-resolution section}).  The
      closed-loop network-bias validation suppresses the wide-lane
      integers, so AR runs on the float state alone.
\item \textbf{State extraction and trace} (\texttt{writeEpochTraces}).
      The posterior is synced back to the PPP state and the trace rows
      are written.  The optional
      ambiguity-level check (card ALC), the per-satellite consumer
      readout (card CST) and the periodic filter-state snapshot follow.
\end{enumerate}

\paragraph{Post-processing stages skipped in real time.}
Non-causal and batch stages do not fit a streaming pipeline without
look-ahead and are absent:
\begin{itemize}[leftmargin=*, itemsep=0pt]
\item the backward RTS \dcref{sec:kf}{smoother} is disabled under
      \texttt{AppMode::RealTime}, because smoothing requires the full
      forward trajectory and cannot run causally;
\item geometry relinearisation (the IEKF second pass) is a
      post-processing-only refinement and is not replicated;
\item precise SP3/CLK products are not loaded when a live
      broadcast$+$SSR source replaces them (a real-time \emph{file}
      replay without corrections still loads them).  Orbit and clock are
      gated separately, so a configuration with precise orbits and
      network clocks loads the SP3 file and skips the CLK file;
\item the stages that need a GIM have no source in the stream:
      higher-order ionosphere, GIM ionosphere weighting and adaptive
      ionospheric process noise are skipped with a warning when enabled.
\end{itemize}

\subsection{Assessing a real-time run}\label{sec:rt-assess}

A streaming run has no answer key, unlike a simulated one.  The
self-referenced statistic of the ordinary state viewer (scatter about
the run's own median) cannot see an error \emph{common} to every epoch,
and a correction-driven pipeline is most likely to make that kind of
error.  A static receiver on a surveyed mark supplies the
missing datum.  Let
$\rxpos_{\text{ref}}$ be the surveyed antenna position
(\texttt{receiver.reference\_position\_ecef}) with geodetic latitude and
longitude $(\varphi,\lambda)$, and $\hat{\rxpos}_k$ the filter's
posterior position at epoch $k$.  The error is the ECEF difference
rotated into the local horizon frame at the \emph{reference} point,
\begin{equation}
\begin{bmatrix} e_k \\ n_k \\ u_k \end{bmatrix}
= \bm{R}(\varphi,\lambda)\,
  \bigl(\hat{\rxpos}_k - \rxpos_{\text{ref}}\bigr),
\qquad
\bm{R} =
\begin{bmatrix}
-\sin\lambda & \cos\lambda & 0\\
-\sin\varphi\cos\lambda & -\sin\varphi\sin\lambda & \cos\varphi\\
 \cos\varphi\cos\lambda &  \cos\varphi\sin\lambda & \sin\varphi
\end{bmatrix},
\label{eq:rt-enu-error}
\end{equation}
with $\varepsilon^{\text{3D}}_k = \lVert(e_k,n_k,u_k)\rVert$ and
$\varepsilon^{\text{H}}_k = \lVert(e_k,n_k)\rVert$.  The difference is
taken in ECEF before the rotation, so the rotation origin cancels and
\eqref{eq:rt-enu-error} is exact rather than a local-tangent
approximation.

\paragraph{Convergence must be suffix-stable.}
The first epoch whose error falls below a threshold measures the noise
of a solution, not the solution.  A float PPP-RTK trajectory dips under
any loose bound long before it \emph{holds} it, and a run that later
diverges scores the same as one that does not.  The convergence epoch
for threshold $\theta$ is therefore defined as
\begin{equation}
k^{\star}(\theta) = \min\Bigl\{\,k :
  \max_{j \ge k}\ \varepsilon^{\text{3D}}_j \le \theta \,\Bigr\},
\label{eq:rt-converge}
\end{equation}
the first epoch from which the error stays within $\theta$ for the
\emph{entire remainder} of the run ($k^\star = \varnothing$ when no such
epoch exists, the correct verdict for a run that never settles).  The
suffix maximum is computed in one reverse cumulative pass.  The suffix
function counts a non-finite error as a violation, but the scorecard
first drops every trace row whose position is non-finite, so an epoch
without a position is skipped rather than breaking the hold.  The metric
is reported as the ladder $\theta \in \{1.0, 0.5, 0.2, 0.1\}$\,m rather
than a single number, which keeps it informative for runs of very
different quality.  The steady-state statistics are taken over
$k \ge k^{\star}$ of the coarsest achieved rung, so the convergence
transient is excluded from the accuracy claim by construction; when no
rung is held, they fall back to the second half of the run and the
summary flags the fallback.  The window is stated alongside the number
it produced.

\paragraph{Covariance credibility.}
Grading the solution against the surveyed point also tests whether the
filter's own uncertainty is credible.  With $\covM_k^{xx}$ the position
block of the posterior covariance, the rotation-invariant scalar
\begin{equation}
\sigma^{\text{3D}}_k = \sqrt{\operatorname{tr} \covM_k^{xx}}
= \sqrt{\sigma_x^2 + \sigma_y^2 + \sigma_z^2}
\label{eq:rt-sigma3d}
\end{equation}
is compared directly with $\varepsilon^{\text{3D}}_k$.  The trace is
invariant under the rotation \eqref{eq:rt-enu-error}, so no per-axis
horizon $\sigma$ has to be constructed from a trace that carries only
the diagonal.  A float solution whose error persistently exceeds
$3\sigma^{\text{3D}}$ is \emph{overconfident} as well as inaccurate.
This is the usual signature of unmodelled correction error absorbed by
ambiguity and bias states whose formal variances do not reflect it.  The
simulator dashboard applies the same covariance-optimism diagnostic
against its answer key; here the surveyed mark takes the place of the
truth trajectory.

\paragraph{What the reference does not certify.}
The surveyed point grades \emph{position} only.  A static run cannot
separate a receiver-clock error from a common-mode range error, and
under BSSD (\Cref{sec:rt-epoch}) there is no receiver-clock state to
grade.  The per-satellite bias states are gauge-fixed, so their absolute
values cannot be checked against an external reference either.  A
position that grades well therefore certifies the geometry and the
correction stack \emph{jointly}, not any single component.  The
real-time scorecard viewer (\texttt{pyviz/gnssviz/viewers/rt\_scorecard.py})
of the report built by \texttt{pyviz/run\_report.py --run-config}
implements \eqref{eq:rt-enu-error}--\eqref{eq:rt-sigma3d}.  An explicit
\texttt{--ref-ecef} overrides the configured reference position.

\ifSubfilesClassLoaded{\printbibliography}{}
\end{document}
